\documentclass[12pt]{article}
\usepackage[T1]{fontenc}
\usepackage[utf8]{inputenc}
\usepackage{lmodern}
\usepackage{textcomp}
\usepackage{enumitem}
\usepackage{geometry}
\geometry{margin=1in}
\begin{document}
\begin{center}
Answer Key - Physics - Graduate Level Assessment 1 \
\small (Answers are ordered exactly as in the question paper.)
\end{center}

\section*{Multiple Choice}
\begin{enumerate}[label=\arabic*.]
\item \textbf{Answer:} D
\\ \textit{Options:} A) 5.5 \ensuremath{\times} 10^-13 g; B) 2.5 \ensuremath{\times} 10-13 g; C) 2.0 \ensuremath{\times} 10^-12 g; D) 4.7 \ensuremath{\times} 10-13 g; E) 8.6 \ensuremath{\times} 10^-13 g
\\ \textit{Note:} The increase in mass of the water is calculated using Einstein's mass-energy equivalence principle (E=mc²), where the energy supplied (10 calories, equivalent to 41.84 joules) results in an extremely small mass increase of approximately 4.7 × 10^-13 g due to the large value of the speed of light squared in the denominator.
\item \textbf{Answer:} A
\\ \textit{Options:} A) 1750.0; B) 1500.0; C) 2200.0; D) 2100.0; E) 1400.0
\\ \textit{Note:} The correct answer of 1750.0 m/s for the mean molecular speed of hydrogen (H$_{2}$) is derived from the formula v = sqrt(3 kT/m), where k is the Boltzmann constant, T is the absolute temperature, and m is the molecular mass of hydrogen, resulting in a high speed due to hydrogen's low molecular mass.
\item \textbf{Answer:} C
\\ \textit{Options:} A) two-thirds; B) twice; C) the same; D) half; E) four times higher
\\ \textit{Note:} The height of mercury in both barometers remains the same because atmospheric pressure acting on the surface of the mercury is constant, causing the mercury to rise to the same height regardless of the cross-sectional area of the tubes.
\item \textbf{Answer:} A
\\ \textit{Options:} A) Virgo; B) Perseus; C) Ursa Major; D) Sagittarius; E) Aquila
\\ \textit{Note:} Virgo is located in a region of the sky that is not densely populated by stars in the Milky Way, as it lies farther away from the galactic plane compared to other constellations that are directly along the Milky Way.
\item \textbf{Answer:} A
\\ \textit{Options:} A) 48 ft per sec; B) 64 ft per sec; C) 30 ft per sec; D) 72 ft per sec; E) 16 ft per sec
\\ \textit{Note:} The correct answer is 48 ft per sec because the final speed can be calculated using the formula \( v = u + at \), where \( u \) (initial velocity) is 0, \( a \) (acceleration) is 16 ft/s², and \( t \) (time) is 3 seconds, resulting in \( v = 0 + (16 \, \text{ft/s}^2 \times 3 \, \text{s}) = 48 \, \text{ft/s} \).
\end{enumerate}

\section*{True / False}
\begin{enumerate}[label=\arabic*.]
\setcounter{enumi}{5}
\item \textbf{Answer:} False
\\ \textit{Options:} A) True; B) False
\\ \textit{Note:} The angular magnification of a refracting telescope is determined by the focal lengths of the objective and eyepiece lenses, not merely by the separation of the lenses.
\item \textbf{Answer:} True
\\ \textit{Options:} A) True; B) False
\\ \textit{Note:} The rest mass of a particle can be calculated using the relation \(E^2 = p^2 c^2 + m_0^2 c^4\), and substituting \(E = 5.0 \text{ GeV}\) and \(p = 4.9 \text{ GeV}/c\) yields a rest mass of approximately \(1.0 \text{ GeV}/c^2\).
\item \textbf{Answer:} True
\\ \textit{Options:} A) True; B) False
\\ \textit{Note:} The efficiency of a device is calculated by dividing the useful work output (40 J) by the input energy (100 J) and multiplying by 100, resulting in an efficiency of 40\%, which is not 25\%, making the statement false.
\item \textbf{Answer:} False
\\ \textit{Options:} A) True; B) False
\\ \textit{Note:} Doubling the thickness of a neutral density filter does not always halve its transmittance because the transmittance depends on the filter's optical density and the specific material properties, which can vary with thickness and wavelength.
\item \textbf{Answer:} True
\\ \textit{Options:} A) True; B) False
\\ \textit{Note:} Density is defined as mass divided by volume, so if the volume of an object doubles while its mass remains constant, the density must halve, as the same mass is now distributed over a larger volume.
\end{enumerate}

\section*{Long Form Response}
\begin{enumerate}[label=\arabic*.]
\setcounter{enumi}{10}
\item \textbf{Answer:} The shape of the illuminated area on a screen positioned 1 cm away from a tiny plane mirror measuring 2 mm \ensuremath{\times} 2 mm will be circular, while at a distance of 2 m, the shape will transform into a rectangular area. At$_{1}$ cm, the small size of the mirror causes light rays reflecting off its surface to diverge rapidly, producing a circular pattern due to the close proximity of the screen. Conversely, at 2 m, the increased distance allows the reflected rays to converge more uniformly, leading to a rectangular projection that corresponds to the dimensions of the mirror. This transformation illustrates the principle of geometric optics, where the distance from the mirror significantly influences the shape of the light distribution on the screen.
\\ \textit{Note:} correct because the small size of the mirror causes light rays to diverge rapidly at a short distance, creating a circular illuminated area, while the greater distance allows the rays to converge more uniformly, resulting in a rectangular shape.
\item \textbf{Answer:} To produce a virtual image 25 cm from the eye using a converging lens with a focal length of 5.0 cm, the required object distance can be calculated using the lens formula: \( \frac{1}{f} = \frac{1}{d_o} + \frac{1}{d_i} \). Here, \( d_i = -25 \) cm (negative since the image is virtual), leading to \( \frac{1}{5} = \frac{1}{d_o} - \frac{1}{25} \). Solving this equation gives \( d_o \approx 4.2 \) cm. The magnification \( M \) is calculated using the formula \( M = -\frac{d_i}{d_o} \), yielding \( M \approx 5.9 \). These values indicate that the optical system is capable of significantly enlarging the image while allowing it to be positioned within a comfortable viewing distance, emphasizing the practical applications of magnifying lenses in various optical devices.
\\ \textit{Note:} The answer correctly applies the lens formula to determine the object distance needed to achieve a virtual image at the specified distance from the eye, while also accurately calculating the magnification to illustrate the effectiveness of the lens as a simple magnifier.
\end{enumerate}

\vfill
\noindent\textit{End of Answer Key}
\end{document}